\chapter{Software Architecture Design}
\label{chap:software-architecture-design}

\section{System Architecture Overview}
\label{section:architecture-overview}

KUru is a three-tier web application with a distinct AI layer. The frontend (Next.js) communicates with a Python FastAPI backend, which orchestrates grounded retrieval, recommendation, and structured program-data services. The architecture is described in two layers: the evaluated MVP architecture and the target-system architecture.

\subsection{Evaluated MVP Architecture}
\label{subsection:mvp-architecture}

The evaluated MVP uses Supabase PostgreSQL with pgvector as the primary data layer. The backend supports: (1) a structured RAG pipeline for curriculum, TCAS, and fee Q\&A; (2) RIASEC interest-profile storage and scoring; (3) simplified Pipeline~B2 recommendation based on semantic matching against indexed program, PLO, and course content; (4) semantic program search; and (5) static PLO/profile visualisation. Gemini 2.5 Flash Lite via OpenRouter is used for grounded answer generation and explanation, while local multilingual E5 embeddings are used for Thai-English retrieval.

Neo4j, O*NET career matching, behavioural re-ranking, portfolio checking, program comparison, and curriculum timeline generation are not required for MVP pass/fail evaluation. They remain part of the target-system design.

\subsection{Target-System Architecture}
\label{subsection:target-architecture}

The target system extends the MVP with a Neo4j knowledge graph connecting faculties, PLOs, skill clusters, and careers; O*NET occupation matching; behavioural personalisation; portfolio readiness analysis; curriculum timeline extraction; and multi-program comparison. In that target architecture, both RAG and recommendation read from a shared data layer consisting of Supabase (PostgreSQL + pgvector for embeddings) and Neo4j (knowledge graph on Railway).

\begin{figure}[p]
    \centering
    \includegraphics[width=\textwidth]{diagrams/system-architecture.png}
    \caption{KUru System Architecture}
    \label{fig:system-architecture}
\end{figure}

Figure~\ref{fig:domain-model} shows the target-system domain model: the eleven core entities and their relationships. The MVP implements the subset centred on \textbf{Student}, \textbf{RIASECProfile}, \textbf{Program}, \textbf{Faculty}, \textbf{PLO}, \textbf{TCASRecord}, and \textbf{Course} stored in Supabase. Phase~2 adds persistent \textbf{InteractionLog}, \textbf{PortfolioCriteria}, \textbf{SkillCluster}, and \textbf{Career} relationships through Neo4j. In the target system, PLOs are linked to \textbf{SkillCluster} nodes in Neo4j (the \textit{develops} relationship), and \textbf{Career} nodes are linked to the same clusters (the \textit{requires} relationship) --- this shared vocabulary enables the multi-hop recommendation path from student interests to program outcomes.

\begin{figure}[htbp]
    \centering
    \includegraphics[width=\textwidth]{diagrams/domain-model.png}
    \caption{KUru Domain Model}
    \label{fig:domain-model}
\end{figure}

\section{Software Design}
\label{section:software-design}

\subsection{RAG Pipeline Sequence}
\label{subsection:rag-pipeline}

When a user submits a query to the curriculum chatbot: (1) the query is embedded locally using \texttt{intfloat/multilingual-e5-base}; (2) the embedding is used to perform a cosine similarity search against the pgvector store of มคอ.2 document chunks; (3) the top-$k$ retrieved chunks are assembled into a prompt context along with the user query and the conversation history; (4) Gemini 2.5 Flash Lite via OpenRouter generates a grounded response in the requested language; (5) source citations are appended to the response. In the current POC, citation chips expose provenance such as source file, section type, and similarity/confidence context, but they do not yet open the exact PDF page or extracted chunk. Deep source inspection is Phase~2; the future API should include \texttt{chunk\_id}, page number, and source URL so the frontend can open a source preview panel or PDF page.

\begin{figure}[p]
    \centering
    \includegraphics[width=\textwidth]{diagrams/rag-pipeline.png}
    \caption{RAG Pipeline Sequence}
    \label{fig:rag-pipeline}
\end{figure}

\subsection{Recommendation Pipeline Sequence}
\label{subsection:recommendation-pipeline}

For the evaluated MVP, recommendation uses the simplified Pipeline~B2 path. RIASEC is not treated as proof that a student will be satisfied with a program after graduation. Instead, it is used as an interest signal that can be converted into curriculum search intent. When a student completes the interest discovery flow, the backend converts the RIASEC profile into interest themes and a semantic query, retrieves semantically similar program/PLO/course chunks from Supabase pgvector, aggregates evidence by program, computes a curriculum-fit score, and generates a cited plain-language explanation from the retrieved evidence. This is the recommendation behaviour evaluated for MVP pass/fail.

For example, a high Investigative--Realistic profile is converted into themes such as analysis, problem solving, systems, technical design, experiments, and hands-on laboratory work. These themes are embedded as a semantic query and matched against curriculum evidence such as PLOs, course descriptions, project-based learning statements, laboratory requirements, and technical skill outcomes. A program ranks highly only when retrieved curriculum chunks support the match.

\begin{figure}[htbp]
    \centering
    \fbox{\parbox{0.92\textwidth}{
    \textbf{MVP recommendation flow:}
    RIASEC answers $\rightarrow$ 6D RIASEC profile $\rightarrow$ top dimensions / Holland code $\rightarrow$ interest themes and semantic query $\rightarrow$ Supabase/pgvector search over program, PLO, and course chunks $\rightarrow$ group evidence by program $\rightarrow$ curriculum-fit score $\rightarrow$ ranked programs with evidence-grounded explanation.
    }}
    \caption{MVP Recommendation Pipeline: RIASEC-to-Curriculum Semantic Matching}
    \label{fig:mvp-recommendation-pipeline}
\end{figure}

The full target-system recommendation pipeline proceeds in five stages using two independent parallel signals:

\begin{enumerate}[leftmargin=80pt]
    \item \textbf{RIASEC vector construction.} The twelve-step adaptive elicitation (UC-01) produces a weighted 6-dimensional RIASEC vector. Steps~1--6 (Likert screens) produce raw scores per dimension (max 20 each). Step~7 applies a global confidence scalar (1.0, 0.75, or 0.5). Step~8 (adaptive pairwise) adjusts scores for dimension pairs with delta $< 3$. Steps~9--11 (scenario questions) adjust scores proportionally per response. Step~12 applies the optional dealbreaker filter, zeroing out explicitly rejected dimensions. The resulting vector is L2-normalised to a unit vector. The MVP uses this vector to form a Pipeline~B2 curriculum search intent; the target system also uses it in Pipeline~A and Neo4j Pipeline~B1.

    \item \textbf{Parallel signal computation.} Two independent pipelines run concurrently:

    \textbf{Pipeline A (Career-side):} The student RIASEC vector is compared against precomputed O*NET occupation RIASEC interest profiles in Supabase pgvector using cosine similarity. O*NET interest profiles are empirically derived from surveys of employed workers, making this step grounded in real-world data. The top 10--15 occupations form the internal career space; the top 7 are surfaced to the student in the Career Explorer (UC-06). An A-score per KU program is computed based on career-program pathway alignment. O*NET's 35-skill taxonomy is used for Career Explorer display in UC-06 only --- it is not used as a bridge to Neo4j in the ranking pipeline.

    \textbf{Pipeline B (Curriculum-side):} Two sub-signals computed independently and combined:
    \begin{itemize}[leftmargin=40pt]
        \item \textbf{B1 --- Neo4j PLO match:} The student's RIASEC vector is converted to a SkillCluster weight vector and matched against KU program PLO SkillCluster profiles via Neo4j graph query. B1-score is the cosine similarity between the student weight vector and each program's PLO skill profile.
        \item \textbf{B2 --- Semantic course content match:} The student's interest query is matched against course description embeddings in Supabase pgvector via semantic search over ingested มคอ.2 course content. B2-score is the weighted mean similarity of the top retrieved course chunks per program.
    \end{itemize}
    Combined B-score $= 0.4 \times \text{B1} + 0.6 \times \text{B2}$.

    The two pipelines are independent --- an error in Pipeline~A does not propagate into Pipeline~B.

    \item \textbf{Score synthesis.} Programs are ranked by a final score combining both pipeline outputs:
    \[\text{score} = 0.35 \times \text{A-score} + 0.65 \times \text{B-score}\]
    The top-$N$ programs are selected for display.

    \item \textbf{Explanation generation.} The top-$N$ ranked programs, together with their Pipeline~A details (matched O*NET occupations) and Pipeline~B details (curriculum themes and matched PLOs), are passed to Gemini 2.5 Flash as structured context. Gemini generates a plain-language explanation for each program referencing both career alignment and curriculum fit. Explanations are generated in Thai or English depending on the student's language setting.

    \item \textbf{Behavioural re-ranking (registered users only).} For registered users with accumulated interaction data, the pipeline scores from Stage~3 are blended with a behavioural fit score derived from the individual student's own interaction history (not from other users --- true collaborative filtering across students is deferred to Phase~2). The blend ratio is controlled by $\alpha$: new users start at $\alpha = 1.0$ (pure pipeline score); $\alpha$ decays toward 0.2 as interaction weight accumulates according to the threshold schedule documented in UC-07 (3--5 meaningful signals $\rightarrow$ $\alpha \approx 0.5$; regular user $\rightarrow$ $\alpha \approx 0.2$). The final ranking score is:
    \[\text{final score} = \alpha \times \text{pipeline score} + (1 - \alpha) \times \text{behavioural fit}\]
    Guest interaction signals are session-scoped and do not update $\alpha$; guests always begin a new session at $\alpha = 1.0$. Fallback hierarchy when behavioural data is sparse: (1) sufficient interaction history $\rightarrow$ normal blended signal; (2) limited interaction history $\rightarrow$ $\alpha$ held higher; (3) no interaction history $\rightarrow$ pure pipeline score ($\alpha = 1.0$). Interaction weights are documented in UC-07.
\end{enumerate}

\begin{figure}[p]
    \centering
    \includegraphics[height=0.9\textheight, keepaspectratio]{diagrams/recommendation-pipeline.png}
    \caption{Phase~2 Target-System Recommendation Pipeline --- Five Stages with Two Parallel Signals}
    \label{fig:recommendation-pipeline}
\end{figure}

\subsection{Interest Elicitation Flow}
\label{subsection:elicitation-flow}

The interest elicitation flow (UC-01) produces the RIASEC vector that seeds the MVP recommendation pipeline. In Phase~2, the same vector can also seed the Career Explorer. The flow is adaptive: the pairwise comparison step (Step~8) is shown only when at least one dimension pair has a score delta below the 3-point ambiguity threshold; students with a clearly differentiated Likert profile skip directly to the scenario questions. Figure~\ref{fig:elicitation-flow} shows the full conditional logic from entry to RIASEC vector output.

\begin{figure}[p]
    \centering
    \includegraphics[height=0.9\textheight, keepaspectratio]{diagrams/elicitation-flow.png}
    \caption{Interest Elicitation Flow (UC-01) — 12-Step Adaptive Process}
    \label{fig:elicitation-flow}
\end{figure}

\section{AI and Data Design}
\label{section:ai-data-design}

\subsection{Data Sources and Collection}
\label{subsection:data-sources}

\begin{itemize}[leftmargin=80pt]
    \item \textbf{KU curriculum documents (มคอ.2):} Provided directly by KU faculty advisors for all KU programs. Thai-language PDFs of varying length (typically 20--80 pages per program). These are the primary knowledge base for both the RAG pipeline and the PLO extraction pipeline. Documents are version-tagged by academic year and re-ingested when updated.

    \item \textbf{O*NET occupational database \cite{onet2024}:} Publicly available from the US Department of Labor. Contains two data layers used by KUru: (a) \textit{Interests} --- the RIASEC profile for each occupation, derived from surveys of employed workers, used for career matching from the student's RIASEC vector; and (b) \textit{Skills} --- 35 standardized skills per occupation with importance and level scores, displayed in the Career Explorer (UC-06) to show students what skills their matched careers require. Not used as a bridge to Neo4j in the ranking pipeline --- Pipeline~B1 queries Neo4j directly from the RIASEC-derived SkillCluster weight vector. Downloaded as structured data (CSV/JSON). Occupation titles are mapped to Thai equivalents for student-facing display; skill and interest data are used directly as the underlying structure is language-independent.

    \item \textbf{TCAS admission data:} Ingested from three sources in priority order. (1)~\textit{KU faculty-provided project PDFs} received via Google Drive --- each PDF (e.g., โครงการช้างเผือก.pdf) may contain admission records for multiple programs across multiple faculties. A batch downloader retrieves all PDFs recursively. Each PDF is classified as tcas / mko2 / portfolio / unknown using rule-based keyword detection on the first page. TCAS PDFs are processed by a structured extraction pipeline using Gemini to parse per-program TCASRecord objects covering project name, round, faculty, program name, quota, GPAX minimum, exam score criteria, portfolio requirements, and application deadlines. (2)~\textit{mytcas.com} --- scraped as a gap-filling layer for programs not covered by faculty PDFs, using the mytcas\_code field in the canonical program registry. (3)~\textit{Faculty portfolio criteria PDFs} --- retained as a Phase~2 source for the Portfolio Readiness Checker rather than an MVP pass/fail requirement. Available sources are resolved against a canonical program registry using fuzzy name matching with PyThaiNLP normalization and sequence similarity scoring (threshold 0.85). Records below this threshold are flagged for manual review. A Supabase materialized view (\texttt{programs\_unified}) joins มคอ.2 structured data, TCAS records, and available admission metadata into a unified program record keyed on program\_id, refreshed after each ingestion run.

    \item \textbf{Interaction log (Phase~2):} Generated by the system as registered users engage with KUru. Stored in Supabase PostgreSQL. Records program saves, card expansions, TCAS guide views, chatbot queries, and card views, each with a timestamp, interaction weight, and associated program identifier. Used to compute the behavioural fit score in Stage~5 of the target-system recommendation pipeline. Not required for MVP pass/fail evaluation.

    \item \textbf{Faculty portfolio criteria documents (Phase~2):} Published by KU faculties for TCAS Round~1 and Round~2 applications. Ingested via the same PDF extraction pipeline as มคอ.2 using Gemini in structured extraction mode to produce a per-faculty, per-round criteria schema. Stored in Supabase PostgreSQL alongside TCAS admission data. Updated at the start of each TCAS cycle. Used by the Portfolio Readiness Checker.

    \item \textbf{Course structure data (partial MVP / Phase~2):} Extracted from มคอ.2 during ingestion as a structured layer stored separately from the embedding chunks where available. MVP uses available course and PLO text for RAG, semantic search, and simplified recommendation. Full per-year course timeline visualisation is Phase~2.
\end{itemize}

มคอ.2 document ingestion uses a six-stage pipeline:

\begin{enumerate}[leftmargin=80pt]
    \item \textbf{PDF classification.} Pages are classified by text density and Thai character density (minimum 20\% Thai Unicode characters for Thai-language documents, minimum 100 characters per page) into born-digital, scanned, pure-image, or mixed types.

    \item \textbf{Tiered text extraction.} The current curriculum ingestion pipeline first runs PyMuPDF on every PDF page. Pages with very low text yield are rendered to images and sent to Typhoon OCR (\texttt{typhoon-ocr}) when \texttt{TYPHOON\_API\_KEY} is available; recent targeted re-ingests mostly record \texttt{pymupdf+typhoon\_pages}. If the whole PDF remains below the scanned-document threshold, the ingest script treats it as a full scanned PDF and calls the isolated \texttt{ocr\_extractor.extract\_with\_ocr()} path. That full scanned-PDF path uses the configured \texttt{OCR\_MODEL}: by default direct Gemini \texttt{gemini-2.5-flash}, or Typhoon/OpenRouter if explicitly configured, with Tesseract as a last local fallback when vision OCR yields too little text. Each page is tagged with its extraction method (\texttt{pymupdf}, \texttt{typhoon\_page}, \texttt{scanned}, or \texttt{gemini\_ocr}) for coverage reporting.

    \item \textbf{Section-aware chunking.} PLO sections, course descriptions, admission requirements, and curriculum structure sections are chunked separately (target 500 tokens, 50-token overlap) to preserve semantic coherence. Each chunk is tagged with section type (\texttt{plo} / \texttt{course} / \texttt{admission} / \texttt{curriculum\_structure} / \texttt{general}).

    \item \textbf{Structured extraction.} After text extraction and chunking, Gemini text mode (\texttt{LLM\_MODEL=gemini-2.5-flash-lite} by default) extracts program metadata, PLOs, course lists, year timeline fields, and curriculum mapping into structured JSON. The MVP stores these fields in Supabase for program detail pages, semantic search, RAG grounding, and coverage reporting. Neo4j population from PLO/skill data is Phase~2.

    \item \textbf{Embedding generation.} The current POC uses the local multilingual E5 embedding model (\texttt{intfloat/multilingual-e5-base}) to embed chunks before upserting them to Supabase/pgvector. This avoids external embedding quota limits. Chunks from failed OCR pages are excluded.

    \item \textbf{Target-system graph population and quality reporting.} For Phase~2, Program $\rightarrow$ PLO $\rightarrow$ SkillCluster nodes and edges are created in Neo4j from the structured PLO extraction output. For the MVP, the corresponding quality report focuses on whether curriculum, TCAS, fee, and PLO chunks are retrievable from Supabase/pgvector and whether documents with OCR or extraction failures are flagged for manual review.
\end{enumerate}

\subsection{Model Design}
\label{subsection:model-design}

The primary answer-generation model for the MVP RAG path is Gemini 2.5 Flash Lite via OpenRouter, selected for Thai/English response quality and cost efficiency. Direct Gemini text-mode calls are used for structured extraction where appropriate, and direct Gemini vision OCR is used only in the isolated full scanned-PDF OCR path when configured. The project avoids LangChain or similar orchestration frameworks to reduce latency and maintain explicit control over prompt structure.

Embeddings use the local \texttt{intfloat/multilingual-e5-base} model, with the query/passage prefixes required by the E5 family. This model produces multilingual semantic representations for Thai and English while avoiding Gemini embedding quota limits.

The Phase~2 Neo4j knowledge graph uses a schema with four node types: Faculty, PLO, SkillCluster, and Career. Edges represent: Faculty -[HAS\_PLO]$\rightarrow$ PLO, PLO -[DEVELOPS]$\rightarrow$ SkillCluster, Career -[REQUIRES]$\rightarrow$ SkillCluster. This schema supports the multi-hop queries required for the target-system recommendation: given a student's SkillCluster profile, find Faculties whose PLOs develop those clusters.

The interest discovery interface uses a 12-step elicitation flow grounded in Holland's RIASEC model \cite{holland1997}. Steps~1--6 present one RIASEC dimension per screen with 4 Likert statements each (5-point scale, max score 20 per dimension). Step~7 applies a global confidence scalar (1.0 / 0.75 / 0.5) to all dimension scores based on a single self-reported confidence question. Step~8 presents adaptive pairwise forced-choice questions targeting only dimension pairs where the score delta is below a 3-point threshold, with a maximum of 4--6 pairs shown. Steps~9--11 present three scenario questions with A--F role options mapped to RIASEC dimensions, adjusting scores proportionally. Step~12 displays the profile summary with an optional dealbreaker filter that zeros out any explicitly rejected dimensions before L2-normalisation. The full question bank is documented in Appendix~\ref{appendix:A}.

The MVP recommendation ranking uses Pipeline~B2 curriculum-side semantic fit only. The Phase~2 target ranking blends the RIASEC-derived fit score with a behavioural fit score using a blending weight $\alpha$ that decays automatically as the student accumulates interaction data: new users start at $\alpha = 1.0$ (pure RIASEC); after 3--5 meaningful interactions $\alpha \approx 0.5$ (balanced blend); regular users converge toward $\alpha \approx 0.2$ (behaviour-dominated). Interaction signals are weighted by type: program save (1.0), TCAS guide viewed (0.8), chatbot query about a program (0.6), card expanded (0.4), card viewed (0.1).

The baseline for comparison is keyword-only \texttt{ilike} search over the same มคอ.2 corpus without embedding-based retrieval or LLM generation. The full RIASEC-to-topic mapping is documented in Appendix~\ref{appendix:A}.

\subsection{Evaluation}
\label{subsection:evaluation}

\begin{itemize}[leftmargin=80pt]
    \item \textbf{Current POC baseline (May~2026):} The current RAG POC is evaluated with custom LLM-as-judge regression tests tracked in MLflow on a 0--3 scale, including TCAS, fee, alias, bilingual, and missing-data honesty checks. The selected production regression run, \texttt{v8\_structured\_tcas\_fees}, records 72.7\% good answers and an average score of 2.055/3.0 on the fee/TCAS regression suite. The headline retrieval benchmark, \texttt{latest\_submission\_headline\_v7\_rerank\_74pct}, records 74\% good answers and an average score of 2.26/3.0. RAGAS was considered but is not the primary POC metric because custom checks better capture Thai/English KU-specific correctness, structured TCAS/fee evidence, and missing-data honesty. This baseline is reported honestly as POC evidence, not as final MVP success.

    \item \textbf{Final MVP RAG quality:} RAGAS framework metrics \cite{es2023ragas} are added as standardized supplementary metrics for UC-04 and UC-08: Faithfulness (proportion of claims in the answer supported by retrieved context), Answer Relevancy (semantic similarity between answer and question), and Context Recall (proportion of relevant information retrieved). Target: Faithfulness $>$ 0.80, Answer Relevancy $>$ 0.75. Custom regression tests remain required for TCAS, fees, aliases, bilingual behaviour, and missing-data honesty because these behaviours are KU-specific.

    \item \textbf{Recommendation quality:} The MVP does not claim to measure whether a student will remain satisfied with a selected program after four years. That outcome is longitudinal and cannot be observed within the project period. Instead, recommendation quality is evaluated as a short-term relevance and explainability task: given a manually curated set of 30--50 student interest profiles, the system must rank programs judged relevant by faculty advisors, senior students, or documented curriculum evidence. Metrics are Mean Reciprocal Rank (MRR) and NDCG@5, with target MRR $>$ 0.60. Explanations are also checked for faithfulness to the retrieved curriculum evidence. The MVP evaluates the simplified Pipeline~B2 recommender; career-path validation, full hybrid Pipeline~A+B, and Neo4j B1 evaluation are Phase~2.

    \item \textbf{Interest discovery and semantic search usability:} UC-01 completion rate target $>$ 80\%, with System Usability Scale (SUS) score from user testing with 10--15 high school students or KU first-year students. UC-09 is evaluated through semantic query checks and user task completion. Target: SUS score $>$ 70 (acceptable).
\end{itemize}
